% !TeX root = ../../Book3.tex
% 纲要标题：### C.4 局部交数、乘数理想与对偶复形
% 纲要位置：工作记录/计划纲要.md:2887

\section{局部交数、乘数理想与对偶复形}
\label{b3:sec:C04}

% 纲要内容（仅作写作计划，尚非教材正文）：
% * 平面曲线有限长度交数与Tor算例；区分不满足有限长度的自交
% * 对数消解存在性采用主理想化定理；乘数理想的解析—代数比较与选择独立性采用准确引用
% * 正规交叉模型中展开消失阶与局部可积性的对应，不以一行完备化说明代替比较定理
% * 复用第22章局部对偶，区分CM的单一次数与非CM商环的多次数Ext
% #### C.4.1 局部交数与 Tor
% #### C.4.2 对数消解与乘数理想
% #### C.4.3 非 Cohen--Macaulay 环的多次数对偶


% 正文以具体模型说明专题；长证明给出概要并指明技术细节的出处。

相交的厚度、函数沿除子的消失阶与自由消解的对偶，分别产生长度、理想和上同调模。它们依赖的对象也不同：交数研究一对曲线，乘数理想研究光滑簇上的理想及其正有理参数，而局部对偶比较局部上同调与 Ext。下面从具体计算说明这些差别。

\subsection{局部交数与 Tor}
\label{b3:subsec:C04:01}

两条平面曲线在一点相遇时，仅记录这个点不够：横截与相切应有不同的交数。局部商环的长度提供最直接的模型。\cref{b3:def:plane-local-intersection-number}给出复平面曲线的定义；这里同样设 \(k\) 为域，\(S=k[\![x,y]\!]\)，曲线由 \(f,g\) 定义，并假定 \(S/(f,g)\) 为有限长度模；在这个平面模型中，把
\[
i_0(f,g)=\ell_S(S/(f,g))
\]
称为原点处的\term{局部交数}[local intersection multiplicity]。

\begin{example}\label{b3:exm:plane-contact-multiplicity}
设 \(k\) 为域，\(S=k[\![x,y]\!]\)。对 \(f=y\)、\(g=y-x^m\)，其中 \(m\ge1\)，有
\[
S/(f,g)\cong k[\![x]\!]/(x^m),\qquad i_0(f,g)=m.
\]
模的滤过由 \(1,x,\ldots,x^{m-1}\) 给出 \(m\) 个长度一因子。另一方面，用 \(0\to S\xrightarrow{y}S\to S/(y)\to0\) 与 \(S/(y-x^m)\) 张量，乘 \(y\) 化为 \(k[\![x]\!]\) 上乘 \(x^m\)，是单射。因此
\[
\operatorname{Tor}_1^S(S/(y),S/(y-x^m))=0,
\]
而零次张量积的长度为 \(m\)。这把交数与本卷的自由消解计算联系起来。
\end{example}

若改为两次同一条曲线 \(y=0\)，零次张量积与一次 Tor 都同构于 \(k[\![x]\!]\)，不再具有有限长度。此时不能写下两个无限长度的形式差来定义交数。一般交理论使用循环、等价关系及更系统的 Tor 交替和，需要额外的有限性和维数假设；本节只证明以上有限长度模型。

\subsection{对数消解与乘数理想}
\label{b3:subsec:C04:02}

另一类信息来自函数沿除子的消失阶。为说明乘数理想，先固定光滑不可约复代数簇 \(X\)、非零凝聚理想层 \(\mathfrak a\) 和有理数 \(c>0\)。选取一个\term{对数消解}[log resolution] \(\mu:Y\to X\)：它是光滑簇之间的固有双有理态射，使 \(\mathfrak a\mathcal O_Y=\mathcal O_Y(-F)\)，且有效除子 \(F\) 与例外除子的支集具有简单正规交叉，即局部由若干坐标超平面的并给出。\(\mathcal O_Y(D)\) 表示局部允许的有理函数层，沿每个素除子的极点阶不超过 \(D\) 的相应系数。相对典范除子 \(K_{Y/X}\) 由 Jacobian 行列式记录顶次微分形式的变化。于是定义\term{乘数理想}[multiplier ideal]
\[
\mathcal J(\mathfrak a^c)
=\mu_*\mathcal O_Y\bigl(K_{Y/X}-\lfloor cF\rfloor\bigr),
\]
其中下取整逐除子系数进行。它首先是在 \(Y\) 上规定函数应消失多少阶，再把这些条件收回 \(X\)。定义见\cite[Lecture 1, Definition 1.3]{Lazarsfeld2010ShortCourse}；所用对数消解的存在性来自特征零主理想化定理，其选择独立性则由双有理变换公式保证。

\begin{theorem}\label{b3:thm:log-resolution-multiplier-independence}
设 \(X\) 为光滑不可约复代数簇，\(\mathfrak a\subseteq\mathcal O_X\) 为非零凝聚理想层，\(c>0\) 为有理数。\(\mathfrak a\) 有对数消解，且
\(\mu_*\mathcal O_Y(K_{Y/X}-\lfloor cF\rfloor)\) 不依赖所选对数消解 \(\mu\)。
\end{theorem}
\begin{proofsketch}
存在性采用特征零的主理想化定理：经过有限次沿光滑中心的吹起，理想的总变换成为简单正规交叉除子的可逆理想，且与例外除子的支集同时具有正规交叉。该结果及理想层版本见\cite[Theorem 1.10, Remark 1.18]{BierstoneMilman1997CanonicalResolution}。构造根据局部不变量选择与已有例外除子正规交叉的光滑中心，吹起后用变换后的不变量继续选择中心，直至理想成为所需的单项式形状。中心选择的相容性与有限终止所用的不变量性质，是\cite[Theorem 1.14, Remark 1.18]{BierstoneMilman1997CanonicalResolution}所建立的深层结论。

先说明所定义的推前层确实为 \(\mathcal O_X\) 的子层。设 \(u\) 为它的一个有理截面。在 \(X\) 任一素除子的泛点附近，固有双有理态射 \(\mu\) 为同构，相对典范除子的系数为零；相应截面条件给出
\(\operatorname{ord}_D(u)\ge\lfloor c\operatorname{ord}_D(\mathfrak a)\rfloor\ge0\)。所以 \(u\) 在 \(X\) 的所有余维一处无极点。光滑 \(X\) 正规，其正则函数环为这些高度一赋值环的交，故 \(u\) 是正则函数。

独立性是\cite[Lecture 1, Proposition 1.6]{Lazarsfeld2010ShortCourse}的结论。用解析方法解释时，要同时采用该文 Proposition 1.8 的解析—代数比较定理：按局部可积性定义的解析乘数理想，是上述代数乘数理想的解析化；其除子形式的双有理变换公式见 Lemma 1.9。下面的坐标计算说明可积条件怎样给出定义中的除子系数。

设 \(f_1,\ldots,f_r\) 为 \(\mathfrak a\) 的局部生成元。对于一个正则函数 \(u\)，考虑
\[
\frac{|u|^2}{(|f_1|^2+\cdots+|f_r|^2)^c}
\]
是否在各点附近可积。更换有限生成集只把分母乘上局部上下有正界的因子，故该条件不变。在任一对数消解的坐标图中，\(f_j\circ\mu\) 的共同因子是 \(z_1^{a_1}\cdots z_l^{a_l}\)，剩余因子共同生成单位理想；Jacobian 的消失阶 \(b_i\) 正是 \(K_{Y/X}\) 的系数。换元后，被积式相当于
\( |u\circ\mu|^2\cdot\prod_i|z_i|^{2(b_i-ca_i)}\)。

一变量积分 \(\int_0^\varepsilon r^{2t+1}\,\mathrm{d}r\) 有限当且仅当 \(t>-1\)。逐变量应用，得到可积性恰好要求
\[
\operatorname{ord}_{z_i=0}(u\circ\mu)
\ge\lfloor ca_i\rfloor-b_i\qquad\text{对所有 }i.
\]
这正是 \(u\) 属于所写推前理想的条件。可积性是在 \(X\) 的解析邻域上定义的，不依赖消解选择；借助上述解析—代数比较与双有理变换定理，便得到代数乘数理想的选择独立性。
\end{proofsketch}

\begin{example}\label{b3:exm:snc-multiplier-ideal}
设 \(c>0\) 为有理数，\(1\le r\le n\)。在 \(X=\mathbb A_{\mathbb C}^n\) 上，设 \(\mathfrak a=(x_1^{a_1}\cdots x_r^{a_r})\)，\(a_i\ge1\)。坐标超平面已为简单正规交叉，恒等映射就是对数消解，且相对典范除子为零，故
\[
\mathcal J(\mathfrak a^c)
=\bigl(x_1^{\lfloor ca_1\rfloor}\cdots x_r^{\lfloor ca_r\rfloor}\bigr).
\]
例如 \(\mathfrak a=(x^2y^3)\)、\(c=1/2\) 时得到 \((xy)\)。这个公式针对一个单项式生成的主理想，不能直接把一个多生成元单项式理想中的各指数分别取整；后者需要考虑整体指数集合。
\end{example}

要让相对典范除子实际参与计算，可取 \(X=\mathbb A_{\mathbb C}^2\)、\(\mathfrak a=(x,y)\)，并吹起原点。第一张图为 \(x=u,y=uv\)，于是
\[
(x,y)\mathcal O_Y=(u),\qquad
\mathrm{d}x\wedge\mathrm{d}y=\mathrm{d}u\wedge(v\,\mathrm{d}u+u\,\mathrm{d}v)=u\,\mathrm{d}u\wedge\mathrm{d}v.
\]
第二张图 \(x=rw,y=w\) 给出同样的消失阶。例外除子 \(E\) 在第一张图由 \(u=0\) 定义，故理想的总变换为 \(\mathcal O_Y(-E)\)，而 Jacobian 说明 \(K_{Y/X}=E\)。这是一个对数消解，因而
\[
\mathcal J((x,y)^c)
=\mu_*\mathcal O_Y\bigl((1-\lfloor c\rfloor)E\bigr)
=(x,y)^{\max\{\lfloor c\rfloor-1,0\}}.
\]
最后一式可按消失阶核验：一个底空间正则函数 \(g\) 拉回后沿 \(E\) 的阶数，就是它在原点的最低总次数，因此要求 \(\operatorname{ord}_E(g\circ\mu)\ge\lfloor c\rfloor-1\)。当右侧非正时，对正则函数没有新增条件；推前截面在底空间正则已由上面的余维一论证保证。例如 \(c=2\) 时理想为 \((x,y)\)，而非 \((x,y)^2\)。这个一次的差正来自顶次微分形式中的 Jacobian 因子：理想阶数与相对典范阶数必须相减。

\subsection{非 Cohen--Macaulay 环的多次数对偶}
\label{b3:subsec:C04:03}

\chapref{b3:ch:22}的局部对偶给出另一种信息。给定 \(n\) 维完备正则局部环 \(S\) 与非零商 \(R=S/I\)，\cref{b3:thm:regular-local-duality}已经给出
\[
D_S\bigl(H_{\mathfrak m}^i(R)\bigr)
\cong\operatorname{Ext}_S^{n-i}(R,S).
\]
若 \(R\) 为维数 \(d\) 的 CM 环，右侧仅在 \(n-d\) 次非零，典范模便足以承载信息。非 CM 时可能有多个非零次数；\cref{b3:def:regular-cover-dual-complex}用整个对偶复形把它们保存下来。低次局部上同调需要相应次数的 Ext 与之对偶，整个复形还保留这些次数之间的联系。这里的模型以给定正则局部覆盖为基础；一般环上对偶复形的存在性属于更进一步的对偶理论。
